\documentclass[12pt,a4paper,oneside]{article}
\usepackage[brazil]{babel}
\usepackage[utf8]{inputenc}
\usepackage[dvipsnames]{xcolor}
\usepackage{listings}
\usepackage{wrapfig}
\usepackage{olymp}

\setcounter{problem}{1}

\begin{document}

\begin{problem}{Idade atual}{idadeatual}{2}

\begingroup
\setlength{\intextsep}{0pt}%
%\setlength{\columnsep}{0pt}

Estamos em 2026, ano do centenário da UFV!

\begin{center}
\includegraphics[scale=0.75]{figuras/BronzeLogotipo100anos.png}
\end{center}

O aniversário da UFV é em agosto, então ela ainda tem 99 anos, já que ainda não fez aniversário neste ano. A partir de 28 de agosto, terá 100 anos!

O curso de Ciência da Computação da UFV também festeja uma data histórica neste ano. A primeira turma do curso foi em 1986, assim o curso está completando 40 anos.

E os estudantes do curso, qual sua idade atual? Faça um programa que calcula essa idade.

% \Notes

\InputFile

A entrada contém apenas uma linha, com um valor inteiro $A$ informando o ano de nascimento do estudante e um caractere informando se já fez aniversário este ano ou não (respectivamente `\texttt{S}' ou `\texttt{N}').

Restrições: $1926 \leq A \leq 2025$.

\endgroup

\OutputFile

Escreva uma linha contendo a idade atual do estudante.

% \Notes

\Example

\begin{example}
\exmp{
2008 S
}{
18
}%
\end{example}

\begin{example}
\exmp{
2008 N
}{
17
}%
\end{example}

\begin{example}
\exmp{
1926 N
}{
99
}%
\end{example}

\begin{example}
\exmp{
2000 S
}{
26
}%
\end{example}

%Comentários sobre o exemplo, se necessário.

\end{problem}

\end{document}